% 第8章 边界、反例与失效条件
\section{边界、反例与失效条件}\label{sec:limits}

前面七章一直在建工具，这一章反过来用工具自查。四组材料——一次被修正的前视错误、一次合成与真实的分歧、一条两读的收益源、以及四种机制性的失效情形——共同划出本文结论的边界。把这些摊开写，比在结论里加一段“存在局限”的免责声明有用。

\subsection{时点安全：纪律也适用于研究者本人}

判据要求策略只用当期可得的信息。这条纪律常被用来检查信号，却容易被放过一样东西——\textbf{分析者自己的筛选动作}。一个真实案例：检验非流动性溢价是否“只在最不流动的一端”，做法是逐级剔掉偏不流动的个股、在剩下的较流动子集里重建多空。第一版用全样本 ADV 中位数做剔除——这本身就是前视：一只股票“全窗口平均是否流动”，包含了窗口末端的信息。改成时点安全的构造（滚动 250 日中位数、每 21 日重算一次并前向填充）后，同一检验给出相反方向的结论：

\begin{center}
\begin{tabular}{@{}>{\raggedright\arraybackslash}p{3.0cm}>{\raggedright\arraybackslash}p{4.2cm}>{\raggedright\arraybackslash}p{4.2cm}@{}}
\toprule
保留（按 ADV 分位） & 前视筛法 & 时点安全 \\
\midrule
$100\%$（全部） & $1.221$ & $1.221$ \\
$70\%$ & $1.849$ & $1.026$ \\
$50\%$ & $1.931$ & $0.831$ \\
$30\%$ & $1.883$ & $0.438$ \\
\bottomrule
\end{tabular}
\end{center}

前视筛法说“越只做流动的股票、溢价越强”（$1.22$ 升到 $1.93$），时点安全筛法说“越只做流动的股票、溢价越弱”（$1.22$ 降到 $0.44$）。前者是筛选时用到未来信息制造的假象。教训有两层：操作层，筛选、universe 构建、样本切分里的每个“先看看哪个好”都要过时点安全这一关；理论层，\textbf{观测装置的时点维度不豁免研究者}——分析者站在观测链内部，不是外部。

\subsection{合成与真实的分歧：演示不等于证据}

在一个合成面板（$N=61$、纯噪声主导）上，按样本内夏普挑选最优候选，选中的是一条纯噪声组合：滚动前向样本外夏普 $0.477$、缩水夏普 $0.053$、过拟合概率 $0.567$——教科书式地演示了选择偏差如何制造假象。把同一套流程搬到真实数据上，同样按样本内挑选，选中的随机组合样本内夏普 $3.649$、样本外 $2.671$、缩水夏普 $0.9991$、过拟合概率 $0.000$，指向“不是选择偏差的产物”。

两组数字看似矛盾，实则揭示一条方法论：合成面板里“按内挑选会选中噪声”是被设定出来的结果——面板由噪声主导，挑选当然选中噪声；真实数据里 $8$ 个价量特征本身都含信号，随机线性组合只是它们的混合，因此合成结论\emph{没有外推基础}。这条教训在本文自身的尺度扫描上同样适用：扫描画出的曲线是真实数据上的\emph{描述性}证据，它确认曲线形状与理论相容，但它不是样本外验证；正式的样本外结论需要滚动切分，本文未对全部候选系统性执行（见 \ref{sec:selflimit} 小节的边界清单）。

\subsection{非流动性：过了成本关，没过可提取关}

第 \ref{sec:empirical} 章的尺度扫描里，非流动性信号是唯一在全部尺度上低成本后为正的候选（含冲击夏普 $1.01$ 到 $1.13$）。但它也是全库中最需要配注释的一条，因为“扛住成本”不等于“收益可提取”。五项检验的结果如下。

分层：按逐期截面把个股分成十层，层序与下期收益的相关系数 $0.975$，首末层价差日均 $10.8$ 个基点——不是尾部驱动的形态，而是横跨整个横截面的单调梯度。

可投资子集（时点安全）：逐级剔除偏不流动个股后，含冲击夏普单调衰减 $1.02\to0.69\to0.47\to0.08$；在流动性最好那 $30\%$ 的子集里溢价基本消失。\textbf{这份收益长在最不流动的一端}——恰恰是“想卖时卖不掉”风险最高的地方。

流动性压力期：截面非流动性恶化最快的 $25\%$ 交易日，日均溢价 $6.6$ 个基点；改善期 $0.2$ 个基点。方向与“流动性补偿”的解释一致，但差值 $t$ 值只有 $+1.47$，不足以作为强证据。

逐年稳定性：全样本夏普掩盖了不短的不利期——2019 与 2020 两年含冲击口径为负（$-1.47$、$-0.30$，合计约占样本四分之一），2021--2023 在 $1.7$ 附近，2025 年最高 $3.12$，2026 年回到零附近。

规模中性化：三种方法（双重排序、截面回归取残差、规模中性组合）试图把收益与规模因子分开，得到的含冲击夏普落在 $0.607$--$0.611$ 一带，按预设判据（$0.3$ 以下记为分离、$0.3$--$0.7$ 记为无法分离、$0.7$ 以上记为不可分离）落在中间档：\textbf{无法确认它是独立于规模的定价}。且规模中性化本身有成本——年化换手从 $14$ 倍升到 $19$ 倍。

合起来的结论是一句不讨喜的话：非流动性这条线索“过了成本关，没过多可提取性关”。收益的形态（单调梯度、尾部依赖、压力期集中）更像流动性风险的定价，而不是一条等待提取的缝隙——这正是第 \ref{sec:reference} 章分解公式的实践注脚：\textbf{分解能扣掉“看得见、可对冲”的暴露，扣不掉“尚未建模”的风险补偿；两者在数据上的边界，常常比理论上模糊。}（非流动性文献的经典讨论见 Amihud 2002\cite{amihud2002}。）

\subsection{判据失效的四种机制}

\begin{enumerate}
  \item \textbf{衰减速率时变。}窗口上界 $\Delta_{\mathrm{hi}}=(2/\lambda)\ln(\delta_0/c)$ 假定 $\lambda$ 是常数。真实中信息吸收速率随关注度、竞争与制度变化；因子在被公开发表后收益显著衰减\cite{mcleanpontiff2016}，就是 $\lambda$ 抬头的一个宏观表现。$\lambda$ 上升会把窗口向左压缩——原本可行的低频网格逐个失效。
  \item \textbf{参与的内生反馈。}观测装置理论把 $\lambda$ 与 $c$ 当作外生参数，但你的参与会改变它们：规模越大，单笔冲击越高（$c$ 上升）；同类策略拥挤，租金耗散越快（$\lambda$ 上升）。两个变化同向压缩窗口，给出“缝隙的自毁”：\textbf{缝隙的存在性依赖于参与它的规模，这在定义层面就无法回避}。容量约束因此不是一个风险管理附录，而是观测装置的一部分。本文没有对这条反馈做正式建模，留作后续。
  \item \textbf{参照工具不可实施。}第三道核验要求参照系由可交易工具构成。A 股市场的做空限制使多空参照工具在多数个股上不可实施——分解链里的“扣除暴露”在纸面上完成了，在盘面上未必做得到。此时正确的口径是退回到仅多头与可得工具（如指数期货、ETF），承认参照系因此变贫、$\alpha$ 的“可提取性”打折。
  \item \textbf{分辨率不足。}尺度窗口的检验需要覆盖窗口两端的网格。本文数据是日线，$\Delta_{\min}=1$ 日，因此只能检验 $\Delta\ge1$ 的一侧；对衰减极快（$\lambda$ 大）的信号，其真实窗口可能整体落在日内，本文的装置根本看不见——这类信号在本文口径下只会呈现“被成本压死”，而不是“窗口在更细的尺度上”。要检验高频侧，需要分钟级数据与相应的成本模型。
\end{enumerate}

\subsection{本文自身的边界}\label{sec:selflimit}

按第 \ref{sec:criteria} 章的记账精神，本文的结论也要配一份边界清单：

\begin{itemize}
  \item \textbf{方向选择含样本内成分}：尺度扫描中每个特征的交易方向取样本内更优者，这是设计选择而非先验约定；它抬高所有候选的基线表现，且该选择应计入试验账本。
  \item \textbf{全样本口径为主}：扫描曲线为描述性证据，未对每个（候选，尺度）组合做滚动样本外检验；样本外验证已有的部分（随机候选与单特征抽查），覆盖面不等。
  \item \textbf{单市场、单时期、日频}：沪深北 400 只、2019--2026、日线；结论对其它市场与更高分辨率的适用性未验证。
  \item \textbf{数据链有两处外部依赖}：后复权因子取自第三方（本地权息文件加密不可解析），已做 400 只全量交叉核验；universe 依赖本地库清单，仍存在幸存者偏差。
  \item \textbf{成本模型是近似}：平方根冲击、参与率上限 $10\%$、ADV 取窗口内中位数；真实执行还受盘中流动性形态、涨跌停排队等本文未建模的因素影响。
\end{itemize}

\subsection{本章要点}

四组自查材料指向同一个态度：\textbf{这套判据最有用的地方，不是它能证明什么，而是它把“为什么不能证明”分成了可以逐条核对的名目}——前视、外推、两读、时变、内生、不可实施、分辨率。一条缝隙声明越是经得起这份名目的逐条核对，它就越接近“可以下注”而非“可以讲故事”。